% Separate post-CP34 candidate, following the independently proved B94.
\subsection{The critical coefficient removes the tracefree first layer}

Keep all field, actual comparison, normalization and break conventions
of B94. In particular $p=5$, $e_K=20$, $\mu _5\subset K$, the residue
field is algebraically closed, the actual good field $L/K$ is tame
of degree24, and we consider different finite special characters.
This subsection concerns only the still hypothetical layer shape
$(\dim W_1,\dim W_2)=(3,4)$ in the both-large residual-wild case.
By B94 its four remaining unordered parameter pairs, written as
$(B_1/5,B_2/5)$, are
\[
                 (7,15),\ (9,13),\ (9,17),\ (11,15). \tag{B-RD95.1}
\]
The original actual $25$-division field and the actual $25$-division
field of the high quotients are both available from the original
$125$-comparison. For each pair, orient these two fields so that
the \emph{critical} parameter pair is $(a,15)$, with $a=7$ or11,
and the opposite pair is $(9,24-a)$. This merely names two existing
$25$-fields; it does not supply $125$-torsion of the high quotients.

\begin{lemma}[A bound at the critical first break]
\label{lem:B-RD95-critical}
\status{Candidate; complete independent audits pending.}
For the critical pair $(a,15)$, let $t<z$ be the lower breaks of
its actual degree-25 point field over its residual field. Let $C$
be the conductor of the off-diagonal character on that residual
field. Then
\[
       t=24a,\qquad C=\frac{z+4t}{5},\qquad
       24a<C\le1086,\qquad C\equiv6\pmod {24}.        \tag{B-RD95.2}
\]
In particular this bound includes the case in which the first
ordinary coefficient at the critical first break vanishes.
\end{lemma}
\begin{proof}
B94's actual normality argument puts the last line at the common
high-translation line. The two coordinate rows of the point field
are exactly
\[
\begin{array}{c|r|r|r}
 m&A&D_0&D_1\\\hline
 a&5a&125a&15000-625a\\
 15&75&1875&5625.
\end{array}                                           \tag{B-RD95.3}
\]
Its absolute index is60000, so every displacement displayed here
is below the fixed-coefficient mixed-characteristic cutoff.
If $t<24a$, the prime-parameter strict row proved in B94 gives
$z=15000-630a-4t>15000-726a$. If $t>24a$, the other strict row
gives $z=15000-750a+t>15000-726a$. For both $a=7,11$, these
bounds exceed5550. But the other coordinate, of value75, has
high displacement5625, and B60 gives $z<5625-75=5550$.
Both strict cases are impossible. Therefore $t=24a$.

The last upper break is $C=(z+4t)/5$, and strict normality gives
$C>t$. If $z\le1110$, then
$C\le(1110+4t)/5<1109$. Suppose $z>1110$. The degree75 term
has five-adic order two; below the cutoff, its first possible
change is at degree $75+25z$. Every other exponent divisible by
five is at least80 and has first possible degree-five change no
earlier than $80+5z$. All these values exceed5625. Thus the first
nonzero ordinary exponent is necessarily
\[
                         J=5625-z.                  \tag{B-RD95.4}
\]
At the first break $t$, the leading degree75 term changes later
than1875, since $75+25t>1875$. Every other five-divisible exponent
capable of an earlier degree-five change is at least80, so its
change has valuation at least $80+5t$. For $a=7,11$,
\[
                            80+5t<1875.              \tag{B-RD95.5}
\]
If $J+t<80+5t$, the first ordinary term would be uniquely first
and would change strictly before the required displacement1875.
There is no possible cancellation: all higher ordinary terms
change later, and every divisible term has just been bounded.
Consequently $J+t\ge80+5t$, so
\[
             z\le5545-4t,\qquad C\le1109.            \tag{B-RD95.6}
\]
No nonzero coefficient at degree $5a+4t$ was assumed; the argument
works whether that critical coefficient is present or absent.

The actual off-diagonal cyclic character has tame weight
$\zeta^6$, as $a\equiv3\pmod4$. Its degree-five field over the
residual field has degree25 over $L$, and hence a tame
uniformizer eigenvalue $\zeta$ (because $25\equiv1\pmod{24}$).
The same injective ramification leading map used in B94 gives
$C\equiv6\pmod{24}$. Thus $C\le1109$ sharpens to $C\le1086$.
These are bounds on the actual cyclic character; the argument
has not asserted existence for any value in that interval.
\end{proof}

\begin{proposition}[The actual field profiles and terminal restrictions]
\label{prop:B-RD95-profiles}
\status{Candidate; complete independent audits pending.}
Let $r,r'$ be the residual breaks in the critical and opposite
orientations, respectively. Then the necessary possibilities are
\[
\begin{array}{c|c|c|c|c}
 a&r&r'&\text{critical diagonal}&\text{opposite off-diagonal}\\\hline
 7&30,54,78&r-12&168&858\\
 11&6,30,54,78,102,126&r+12&264&1362.
\end{array}                                           \tag{B-RD95.7}
\]
The opposite diagonal conductor is216 in both rows. Write $C'$ for
the last column and set
\[
 d=\frac{24a+4r}{5},\qquad
 \beta=\frac{C+4r}{5},\qquad
 \gamma=\frac{C'+4r'}{5}.                            \tag{B-RD95.8}
\]
For the actual compositum $E=F_2F_Q$, its five upper breaks, in
increasing order, are
\[
       \min(r,r'),\quad\max(r,r'),\quad d,\quad
                     \min(\beta,\gamma),\quad\max(\beta,\gamma).
                                                        \tag{B-RD95.9}
\]
They are distinct and all are below576. Put
\[
 V=\psi_{E/L}(\beta),\quad W=\psi_{E/L}(\gamma),\quad
 R=\psi_{E/L}(r),\quad R'=\psi_{E/L}(r'),\quad
 T=\psi_{E/L}(576).                                 \tag{B-RD95.10}
\]
Here $V,W$ are respectively the lower breaks of the critical and
opposite off-diagonal kernel lines. On the actual terminal group
$\operatorname{Gal}(F/E)$, the original diagonal characters obey
\[
\begin{array}{c|c|c}
 \text{original pair}&\operatorname{cond}_E(a_1)\le&
                         \operatorname{cond}_E(d_1)\le\\\hline
 (a,15)&702000-4V&78000a-4W-20R'\\
 (9,24-a)&78000a-4W&702000-4V-20R.
\end{array}                                           \tag{B-RD95.11}
\]
The exact conductor of their sum, the actual restricted cyclotomic
character, is $T$.
\end{proposition}
\begin{proof}
The opposite pair $(9,24-a)$ has both parameters prime to five.
B94's common-high calculation gives its diagonal conductor216 and
its off conductor $3000-126(24-a)$, namely858 or1362.
The critical diagonal is $24a$ by the preceding lemma. B94's
shared-diagonal restriction gives
\[
                         216=24a+4(r-r').             \tag{B-RD95.12}
\]
The critical residual break is6 modulo24 and the opposite one
is18 modulo24. The strict bounds
$r<600-30\cdot15=150$ and $r'<600-30(24-a)$ now give exactly
the lists in (B-RD95.7); the normality inequalities are satisfied
there. They are necessary lists, without a realization claim.

The two original numerical field computations of B94 apply with
these conductor values. The critical $25$-field over its residual
field has conductors
\[
 \psi_{R/L}(r'),\quad24a,\quad C,
\]
where the first expression denotes the actual residual restriction;
the opposite field has the corresponding opposite root conductor,
216 and $C'$. The strict diagonal ordering and the shared diagonal
relation give the common first three upper breaks in (B-RD95.9).
Their last breaks are respectively $\beta,\gamma$.
For $a=7$, equality of the last two would require $C=810$, which
is not6 modulo24. For $a=11$, one has $C\le1086<1410$, so
$\beta<\gamma$. Thus the breaks are distinct. The same exact
kernel-line argument as in B94 shows that the actual compositum
filtration is their union: between the smaller and larger last
breaks, only the kernel line corresponding to the larger break
remains; above the larger break the group is trivial. Below the
smaller last break both independent kernel lines are present, and the common
degree-$5^3$ quotient determines every earlier order. In particular
there are no unobserved intermediate drops.

For explicit computations, the first three lower breaks are
\[
\begin{array}{c|ccc}
 a&\ell_1&\ell_2&\ell_3\\\hline
 7&r-12&r+48&888-4r\\
 11&r&r+60&1080-4r.
\end{array}                                           \tag{B-RD95.13}
\]
Integrating the next slopes125 and625 gives
\[
\begin{array}{c|c|c|c}
 \text{case}&V&W&T\\\hline
 a=7,\ C<810&25C-3312-4r&97938-100C-4r&1391688-100C-2504r\\
 a=7,\ C>810&125C-84312-4r&16938-4r&1715688-500C-2504r\\
 a=11&25C-5520-4r&170730-100C-4r&1089480-100C-2504r.
\end{array}                                           \tag{B-RD95.14}
\]
The largest upper break is at most $1914/5<576$. By B94 the
actual cyclotomic extension has break576 over $L$, so its
restriction to the actual terminal group has exact conductor $T$.

It remains to justify (B-RD95.11) using the original $125$-module
in each orientation. Use precisely the actual point fields and
stabilizers of Lemma~\ref{lem:B-RD94-terminal}. If the original
parameters are $(a,15)$, their minimum low parameter is $a$ and
minimum high-inverse parameter is9; their upper-root lower break
is $V$, lower-root lower break is $W$, and lower-root residual
break is $R'$. For the original pair $(9,24-a)$ these are instead
$9,a,W,V,R$. This is a change of which original module supplies
the primitive $125$ points, not an enlargement of the available
quotient torsion.

Here is the general conversion in these actual fields. If
$V_o,W_o$ denote the original upper and lower kernel lower breaks,
then $E/S_h$ is elementary abelian of degree25 with those two
lower breaks. Its lower-root cyclic quotient has conductor
\[
 z_h=\begin{cases}
 W_o,&W_o<V_o,\\
 (W_o+4V_o)/5,&W_o>V_o.
 \end{cases}                                         \tag{B-RD95.15}
\]
The actual high primitive-point argument bounds the $a_1$
conductor over $S_h$ by $(15600j_{\min}+4z_h)/5$. When this
number exceeds the last upper break of $E/S_h$, its restriction
bound simplifies, in either ordering, to
\[
                 \operatorname{cond}_E(a_1)
                          \le78000j_{\min}-4V_o.     \tag{B-RD95.16}
\]
Similarly, the actual low primitive-point field gives conductor
at most $3120b_{\min}$ over $E_\ell$. The actual cyclic25 field
$E/E_\ell$ has lower breaks $R_o,W_o$, where $R_o$ is the
lower-root residual break. Above its last upper break, restriction
simplifies to
\[
                 \operatorname{cond}_E(d_1)
                    \le78000b_{\min}-4W_o-20R_o.     \tag{B-RD95.17}
\]
All characters here are the original matrix characters on
$\operatorname{Gal}(F/E)$.

For completeness the needed above-last-break tests have generous
uniform margins. The threshold which $15600j_{\min}$ must exceed
is $V_o$ when $W_o<V_o$, and $z_h$ when $W_o>V_o$.
The maxima of these thresholds, and of the last upper break of
$E/E_\ell$, are bounded as follows:
\[
\begin{array}{c|r|r|r|r}
 \text{original pair}&\text{high threshold}&15600j_{\min}&
     \text{low threshold}&3120b_{\min}\\\hline
 (7,15)&51318&140400&16098&21840\\
 (9,17)&80418&109200&10326&28080\\
 (11,15)&29706&140400&28794&34320\\
 (9,13)&143706&171600&4326&28080.
\end{array}                                           \tag{B-RD95.18}
\]
These are direct endpoint bounds from (B-RD95.2), (B-RD95.7)
and (B-RD95.14), splitting $a=7$ at810. For example, when
$a=11$, $W>V$ and $(W+4V)/5=29730-4r\le29706$;
in the mirrored orientation the needed threshold is
$W\le170730-100\cdot270-4\cdot6=143706$.
The other entries follow by the same displayed linear formulas.
Every required inequality is strict. Equations (B-RD95.16)--(B-RD95.17)
therefore apply and give (B-RD95.11), without a disjointness
assumption or an identification of abstract cyclic fields.
\end{proof}

\begin{theorem}[The tracefree first layer is impossible in the residual-wild branch]
\label{thm:B-RD95-tracefree}
\status{Candidate; complete independent audits pending.}
No remaining different-special-character actual $125$-comparison
at $p=5,e_K=20$ can have nontrivial residual wild inertia and layer
shape$(3,4)$. Consequently, in the residual-wild branch its actual
wild image must be the full Iwahori group
\[
 \left\{\begin{pmatrix}1+5a&v\\5w&1+5d\end{pmatrix}:
 a,d,w\in\mathbf Z/25,\ v\in\mathbf Z/125\right\},
 \qquad |W|=5^9.                                    \tag{B-RD95.19}
\]
This is a necessary restriction, not an existence statement.
\end{theorem}
\begin{proof}
Only the four pairs (B-RD95.1) remained after B94. On their same
actual terminal group, the determinant is
$1+25(a_1+d_1)$, whose conductor is at most the larger of the
bounds (B-RD95.11). But it is also the actual cyclotomic character
of exact conductor $T$. Subtracting the two bounds from $T$ gives
\[
\begin{array}{c|r|r}
 \text{original pair}&T-\text{bound}(a_1)&T-\text{bound}(d_1)\\\hline
 (7,15)&676440-2520r&1237200-500C-2500r\\
 (9,17)&1237440-500C-2520r&677400-2500r\\
 (11,15)&365400-2520r&915600-500C-2500r\\
 (9,13)&914400-500C-2520r&365400-2500r.
\end{array}                                           \tag{B-RD95.20}
\]
For $a=7$ the same differences result on either side of810.
Every entry is positive using $r\le78$ in the first two rows,
$r\le126$ in the last two, and $C\le1086$. The smallest
$a_1$ margin is47880 and the smallest $d_1$ margin is50400.
This contradiction excludes all four rows.

B91 and B93 had already reduced the residual-wild case to both
large parameters and the two layer shapes$(3,4),(4,4)$. Only
$(4,4)$ remains. Its image lies in (B-RD95.19), whose order is
$5^9$; the actual layer count is also$5^{1+4+4}=5^9$. Thus the
containing group and actual image agree. No parameter row or
ramification profile has been treated as a realization.
\end{proof}

\begin{remark}[Scope and retained exact supplement]
\label{rem:B-RD95-scope}
\status{Candidate; complete independent audits pending.}
The full first-layer residual-wild Iwahori case is not excluded or
constructed here. The residual-irreducible cases of wild orders
$5^5$ through $5^8$ also remain. The unrestricted first-wild least
index at $p=5,e_K=20$ consequently remains in $\{2,3\}$.
The proof uses the complete independently audited B94 actual-field
and conductor argument, with its full B60/B69/B86/B91--B93
dependencies. No fresh external-source opening is claimed.

The retained standalone exact program
\path{code/reduction_descent/verify_remaining_tracefree_rows.py}
checks all 654 necessary conductor profiles against independent
rational Herbrand integration, the eight closed strict-gap formulas,
all four relative-field threshold maxima, 1,736 critical-bound
inequalities and 26 exact binomial coefficients. Its author run
passed at 04:39:29.825537 UTC on 20 November 2026 in
0.062296249001519755 seconds. The earlier isolated exploratory
probe and its complete output are also retained; it passed in
0.056594914000015706 seconds. No executed B95 computation failed.
The complete source, output, stdout, pre-certificate proof and exact
diff remain retained. These profiles are necessary arithmetic
possibilities only. The proof above supplies the critical coefficient
bound, actual field identification and symbolic strict inequalities
independently of enumeration.
\end{remark}
